%===============================================================================
\section{Related Work}\label{sec:related}
%===============================================================================

\paragraph{Inference and aggregation in data release}
That individually innocuous data items can jointly reveal protected information is an old
observation in database security, where it is known as the inference and aggregation
problem~\cite{lunt1989aggregation,farkas2002inference}; the privacy literature supplies the
best-known illustration, three demographic attributes that jointly single out most
individuals~\cite{sweeney2000simple}. In power systems, public market data have been shown to
expose non-public physical information, for example grid topology from locational marginal
prices~\cite{kekatos2014topology}. Privacy mechanisms for metering data~\cite{mcdaniel2009smartgrid,giaconi2018smartmeter}
and formal models such as differential privacy~\cite{dwork2014dp} or
$k$-anonymity~\cite{sweeney2002kanonymity} bound what can be learned about an individual
record after a release mechanism is applied. They answer a different question from the one
faced by a market operator, who must decide which aggregate operational \emph{fields} may be
published at all.

\paragraph{Data sensitivity assessment under inference leakage}
Data classification and grading is performed field by field in practice, and the Chinese
national standard explicitly asks that aggregation and processing effects be
considered~\cite{gbt43697}, without prescribing how. The closest work to ours is IGNN~\cite{hu2026ignn},
which introduced data inference leakage into sensitivity assessment for cyber-physical power
grids: fields are nodes of an inference graph whose edges encode physical or statistical
inference relations, and sensitivity is propagated along these edges by a graph attention
network from an intrinsic (expert-assigned) score. This established that the risk of a field
depends on its inference relations to sensitive data. Two questions remain open in this line
of work. First, the risk score is calibrated against expert labels, while its ultimate
validation is an attack experiment in which a model is trained on a group of fields; the
attack outcome itself is not the definition. Second, an edge connects two fields, whereas
inference frequently requires several fields jointly, with none of them informative alone.
We take the attack outcome on a field \emph{set} as the definition of risk and derive the
field-level quantity from it.

\paragraph{Feature importance}
Turning set-level values into per-feature scores is the subject of feature attribution.
Shapley-based measures~\cite{shapley1953value,lundberg2017shap,covert2020sage,williamson2020spvim}
average a feature's marginal contribution over all contexts, providing an allocation of predictive value that can dilute the credit of redundant features~\cite{kumar2020shapley,catav2021mci}.
Leave-one-covariate-out (LOCO)~\cite{lei2018loco} and permutation importance~\cite{breiman2001rf,fisher2019mcr,hooker2021unrestricted}
evaluate a single context, the full feature set, so a feature with a near-substitute receives
almost no importance. The Marginal Contribution Index (MCI)~\cite{catav2021mci} takes the
\emph{maximum} marginal contribution over contexts, is characterized axiomatically, and is
immune to dilution. Worst-case assessment also motivates min-entropy and gain-function approaches in quantitative information flow~\cite{smith2009qif,alvim2012gleakage}: an average over contexts is not a guarantee.
Our field-level risk is a budget-constrained MCI evaluated relative to already-public data.

\paragraph{Fast estimation of many refitted models}
Computing the performance of a model refitted on many feature subsets is expensive, which has
motivated lazy~\cite{gao2022lazyvi} and warm-started~\cite{sun2025warmstart} estimators of
variable importance, amortized explanation models~\cite{jethani2022fastshap,covert2023vitshapley},
and conditional generative models trained under random masks~\cite{ivanov2019vaeac,li2020acflow,yoon2020vime,he2022mae}.
Closed-form ridge heads on frozen features are the standard fast-adaptation device in
meta-learning~\cite{bertinetto2019r2d2,lee2019metaoptnet}, and the benefit of sharing a
representation across tasks is well understood~\cite{caruana1997multitask,baxter2000inductive,maurer2016benefit}.
We combine these ingredients into an estimator of set-level inferability and use it inside a scan--certify loop to recover retrained lower bounds on the largest field contributions.

%===============================================================================
\section{From Set-Level Inferability to Field-Level Risk}\label{sec:problem}
%===============================================================================

\subsection{Parties, field roles and threat model}
The \textbf{auditor} is the data owner. It holds the complete history of all operational
fields, including those it regards as sensitive, and grades fields offline before release.
The \textbf{adversary} obtains published fields and is assumed to hold historical auxiliary
labels of the sensitive target, from which it trains a supervised inference model; risk is
therefore measured as out-of-sample \emph{supervised} inferability. The adversary may combine
the field under review with other fields it already holds.

Field roles are fixed by disclosure rules, not by data analysis. Let $G$ be the candidate
fields that the rules allow to be published, $B\subseteq G$ a \emph{public baseline} that is
already released ($B=\emptyset$ unless stated otherwise), $H=G\setminus B$ the fields under
review, and $Y_y$ a sensitive target. Section~\ref{sec:data} instantiates these roles for
four data sets from the categories and release timings of the Chinese electricity-market
disclosure rules~\cite{nea2024disclosure}. An ex-ante estimate of the \emph{same physical
quantity} as a target (e.g.\ a load forecast when the target is metered load) is removed from
$G$, to focus the audit on inference across distinct operational quantities.

\subsection{Set-level inferability}
\begin{definition}[Population inferability]
\[
\Vy^{*}(S) = 1 - \inf_{f}\,\mathbb{E}\big[(Y_y - f(X_S))^{2}\big] \big/ \mathrm{Var}(Y_y),
\]
the infimum taken over measurable $f$. It satisfies $\Vy^{*}(\emptyset)=0$, $0 \le \Vy^{*} \le 1$
and is monotone in $S$.
\end{definition}

\begin{definition}[Empirical inferability]\label{def:vret}
The data are split into training and test rows, with a validation part inside training. For
an attacker family $\mathcal{H}$, every $h\in\mathcal{H}$ is trained on $X_S$ only; the
attacker and its training epoch are selected on validation, and the test $\Rtwo$ is reported:
$\Vy^{\mathrm{ret}}(S) = \Rtwo_{\mathrm{test}}(\hat h_S)$ with
$\hat h_S = \argmax_{h \in \mathcal{H},\,\mathrm{epoch}} \Rtwo_{\mathrm{val}}(h;S)$.
\end{definition}

Finite samples, optimization and selection noise make $\Vy^{\mathrm{ret}}$ slightly
non-monotone. An adversary holding $S$ can train on any subset of $S$ and keep the best, so
the operational evaluation allows subset selection
\begin{equation}\label{eq:closure}
\Vy(S) = \Vy^{\mathrm{ret}}\big(\hat T_S\big), \qquad
\hat T_S = \argmax_{T \subseteq S} \Rtwo_{\mathrm{val}}(T),
\end{equation}
where the subset is chosen on validation and the value is read on test. We write $\Vy$ for the resulting retrained reference throughout. It supplies the set values against which field-risk estimates and release decisions are evaluated.

\subsection{Budgeted marginal risk}\label{sec:mk}
A release audit must account for joint inference, redundant substitutes, and the information already available to an adversary. We therefore evaluate each field in its most informative admissible background, with the background size bounded by $K$. Appendix~\ref{app:requirements} compares these requirements with existing field scores.

\begin{definition}[Budgeted marginal inference risk]\label{def:mk}
For target $y$, public baseline $B$, fields under review $H=G\setminus B$ and budget $K$,
\begin{equation}\label{eq:mk}
\begin{aligned}
\Mik &= \max_{T \subseteq H \setminus \{i\},\, |T| \le K} \Delta_i(B\cup T),\\
\Delta_i(A)&=\Vy(A \cup \{i\}) - \Vy(A).
\end{aligned}
\end{equation}
Any maximizer $T^{*}$ is a \emph{witness background}.
\end{definition}
$M^{(0)}_{i\to y}=\Vy(B\cup\{i\})-\Vy(B)$ is single-field risk, the quantity current practice
grades on; the \emph{escalation} $\Mik-M^{(0)}_{i\to y}\ge 0$ is the risk that appears only in
combination; the \emph{gain matrix} $\Vy(\{i,j\})-\Vy(\{j\})$ lists the increment of $i$ over
every single background.

\begin{proposition}[Relation to MCI]\label{prop:mci}
For monotone $\Vy$, let $v_B(T) := \Vy(B \cup T) - \Vy(B)$ for $T \subseteq H$. Then $v_B$ is monotone,
non-negative and $v_B(\emptyset)=0$; $\Mik = \max_{|T|\le K}[v_B(T\cup\{i\}) - v_B(T)]$ is the
budget-constrained MCI of $v_B$; for $K \ge |H|-1$ and $B=\emptyset$ it is the MCI of $\Vy$.
\end{proposition}
\begin{proof}
Monotonicity of $\Vy$ gives monotonicity of $v_B$, and
$v_B(\emptyset)=0$ by construction; the rest is substitution.
\end{proof}
The axiomatic properties of MCI for monotone evaluation functions~\cite{catav2021mci}---dummy,
symmetry, the self-contribution lower bound $\Mik\ge M^{(0)}_{i\to y}$ and invariance to
duplicated features---therefore transfer to this setting. The finite budget gives the following interpretation for a release audit.

\subsection{From marginal risk to unsafe combinations}
\begin{theorem}[Threshold-free safety margin]\label{thm:margin}
Fix $i$, $y$, $B$, $K$. A number $m\ge 0$ satisfies ``for every threshold $\tau$ and every
background $|T|\le K$: $\Vy(B\cup T)\le\tau-m \Rightarrow \Vy(B\cup T\cup\{i\})\le\tau$''
if and only if $m\ge\Mik$. For a fixed $\tau$, the smallest $m$ with this property can be
strictly smaller than $\Mik$.
\end{theorem}
\begin{proof}
If $m\ge\Mik$ then $\Vy(B\cup T\cup\{i\})=\Vy(B\cup T)+\Delta_i(B\cup T)\le(\tau-m)+m$.
Conversely, for any $T$ choose $\tau=\Vy(B\cup T)+m$; the property gives $\Delta_i(B\cup T)\le
m$. Strictness: with $B=\emptyset$, $K=1$, $\Vy(\{i\})=0.1$, $\Vy(\{j\})=0.5$,
$\Vy(\{i,j\})=0.9$ one has $M^{(1)}_{i\to y}=0.4$, yet at $\tau=0.95$ the necessary margin is $0$.
\end{proof}
Thus $\Mik$ is the answer to ``by how much can publishing $i$ raise the adversary's accuracy,
whatever it already holds within budget''---a property of the field that does not depend on
the threshold later chosen for grading.

\begin{definition}[Criticality and minimal unsafe sets]
Under baseline $B$, $E\subseteq H$ is a $\tau$-\emph{minimal unsafe set} (MUS) if
$\Vy(B\cup E)>\tau$ and $\Vy(B\cup E')\le\tau$ for every $E'\subsetneq E$. Field $i$ is
$K$-\emph{background} $\tau$-\emph{critical} if some $T$ with $|T|\le K$, $i\notin T$, has
$\Vy(B\cup T)\le\tau<\Vy(B\cup T\cup\{i\})$.
\end{definition}

\begin{theorem}[Criticality $\iff$ small MUS]\label{thm:mus}
If $\Vy$ is monotone and $\Vy(B)\le\tau$, field $i$ is $K$-background $\tau$-critical if and
only if it belongs to a MUS of size at most $K+1$.
\end{theorem}
\begin{proof}
($\Leftarrow$) For a MUS $E\ni i$ with $|E|\le K+1$ take $T=E\setminus\{i\}$.
($\Rightarrow$) If $\Vy(B\cup T)\le\tau<\Vy(B\cup T\cup\{i\})$, take an inclusion-minimal unsafe
$E\subseteq T\cup\{i\}$; it is a MUS of size $\le K+1$. If $i\notin E$ then $E\subseteq T$ and
monotonicity gives $\Vy(B\cup T)\ge\Vy(B\cup E)>\tau$, a contradiction.
\end{proof}
Theorem~\ref{thm:mus} makes the field-level flag a faithful projection of the set-level
danger structure: flagging the $\tau$-critical fields misses no field that takes part in a
dangerous combination of at most $K+1$ fields, and a threshold-crossing background identifies an unsafe combination.
The same MUS list gives the smallest set of fields whose withholding breaks every dangerous
small combination (a minimum hitting set), which we use for withholding.


\begin{proposition}[Closed form at $K=1$]\label{prop:k1}
With $I^B_{ij}=\Vy(B\cup\{i,j\})-\Vy(B\cup\{i\})-\Vy(B\cup\{j\})+\Vy(B)$,
$M^{(1)}_{i\to y}=M^{(0)}_{i\to y}+\max\big(0,\max_{j} I^{B}_{ij}\big)$.
\end{proposition}
Hence a field escalates at $K=1$ exactly when it is strictly super-additive with some partner.

\subsection{Background budget}
$\Mik$ is non-decreasing in $K$, so $K$ trades coverage against the cost of evaluating all sets
of size $\le K+1$. We use $K=2$, for three reasons. The threat model is an adversary holding a
handful of other published series, not the whole catalogue. Empirically, high-order
interactions are rare in engineered systems (the hierarchy principle of factorial
experiments~\cite{li2006regularities}), and models restricted to pairwise interactions are
close to unrestricted ones on tabular data~\cite{lou2013ga2m}. Finally, we measure saturation
directly on every data set (Section~\ref{sec:res-robust}). Two caveats apply. There is no general
theorem---a parity function of $K+2$ fields defeats any fixed $K$---and saturation of the marginal gain
does not imply that no larger set crosses the threshold. Section~\ref{sec:res-robust} shows that both
matter in practice: $K=2$ captures most of the marginal risk, but withholding decisions, which must break
every dangerous combination, are better computed at $K=3$.

%===============================================================================
\section{An Amortized Attacker for Set-Level Inferability}\label{sec:model}
%===============================================================================

\subsection{Why an estimator is needed}
Evaluating $\Mik$ for all fields at $K=2$ requires $\Vy$ on every set of size at most three;
with 35 candidates this is 7,175 sets per target, each requiring a family of attackers to be
trained. The table has to be recomputed whenever the public baseline grows, a target is
redefined, or a what-if release is examined. We therefore build one model that answers
$\Vy(S)$ for any $S$ quickly (Fig.~\ref{fig:model}), and use selected retrained backgrounds to assess the scan (Section~\ref{sec:scan}).

\begin{figure*}[!tp]
\centering
\includegraphics[width=\textwidth]{fig3_model.pdf}
\caption{The amortized attacker separates shared representation learning from adaptation to each field set. (a) Random visibility masks train a target-supervised backbone. (b) Each query concatenates masked values, their squares, and frozen hidden features; a separate ridge readout is fitted using training rows, and predictions are averaged across three seeds. The matrices, network nodes, and prediction marks are schematic. All-column reconstruction is evaluated separately as an extension.}
\label{fig:model}
\end{figure*}

\subsection{Masked pre-training}
The adversary does not know in advance which fields it will obtain. We model this by
pre-training one network under random visibility masks. A mask $m\in\{0,1\}^{p}$ is drawn by
first sampling the number of visible fields uniformly and then a uniform subset of that size;
the network $f_\theta(x\odot m,\ m)$ receives the masked values and the mask itself, so that a
hidden field is distinguishable from a zero value. The main estimator predicts the designated sensitive targets from the visible candidates. An extension using \emph{masked reconstruction} predicts every hidden column of the auditor's data
from the visible ones~\cite{yoon2020vime,he2022mae}; it involves no designated target, so one
pre-training serves any later definition of targets and candidates. The last hidden layer
(256 units) is kept as a frozen feature map $\phi_\theta(x\odot m, m)$. Early stopping uses the
validation rows only.

\subsection{Subset-specific closed-form readout}
For a query $(S,y)$ with mask $m_S$ and seed $r$, the feature block is
\begin{equation}
Z_r(S)=[\,x_S,\ x_S^2,\ \phi_{\theta_r}(x_S,m_S)\,],\quad x_S=x\odot m_S,
\end{equation}
standardized with training statistics. A ridge readout
\begin{equation}\label{eq:ridge}
\beta_{r}(S) = \big(Z_r^{\top} Z_r + \lambda I\big)^{-1} Z_r^{\top} y
\end{equation}
is solved in closed form on the training rows, with $\lambda$ chosen by five-fold
cross-validation inside the training rows. The predictions of $R=3$ independently pre-trained
backbones are averaged, $\hat y = R^{-1}\sum_r Z_r^{\mathrm{test}}\beta_r(S)$, and
$\hat V_y(S)=\Rtwo_{\mathrm{test}}(\hat y)$, followed by the monotone closure. Two numerical
safeguards matter on real data: a floor of $10^{-3}$ on feature standard deviations, and
clipping test features to their training range. Without them, units that are almost constant
on the training rows (inactive ReLUs) are amplified on test rows, a few predictions explode and
$\hat V$ collapses to zero for those sets; since $\Mik$ is a maximum over differences, such
collapses turn directly into inflated risk (Section~\ref{sec:res-est}). Pre-training, readout fitting, and the choice of $\lambda$ use training and validation data.

\subsection{Role of the subset-specific readout}
A cheaper alternative lets the network output the prediction directly, $\hat y=f_\theta(x\odot
m_S, m_S)$, and reads $\hat V$ off it. One set of output weights then serves every conditioning pattern, which can limit adaptation to the particular information available in a query. Solving the readout per subset
keeps the expensive representation shared and the combination rule specific.

\subsection{Error propagation}
Write the estimate as $\hat V(S)=V(S)+b+\varepsilon(S)$, where $b$ is a bias common to the sets
of interest and $\varepsilon(S)$ the set-specific part. Since a marginal is a difference of two
nested sets, $\hat\Delta_i(T)-\Delta_i(T)=\varepsilon(T\cup\{i\})-\varepsilon(T)$ and
\begin{equation}\label{eq:prop}
\big|\hat M^{(K)}_{i\to y} - \Mik\big| \;\le\; \max_{|T|\le K}\big|\varepsilon(T\cup\{i\})-\varepsilon(T)\big|,
\end{equation}
independently of $b$. A readout on a frozen representation typically underestimates the fully
trained attacker by a similar amount on nested sets, so a shared component of that error can cancel. Thus bias in $V$ need not translate into equally large error in $M$. This motivates evaluating marginal error alongside set-value error.

\subsection{Scan and certify}\label{sec:scan}
The estimator is used to \emph{scan}: it produces $\hat M^{(K)}$ for every field together with a
ranked list of candidate backgrounds. For \emph{certification},
dedicated attackers are retrained on the top-$k$ backgrounds only, and
$L^{(K)}_{i\to y}=\max_{T\in\text{top-}k}\Delta_i(B\cup T)\le\Mik$ is a lower bound computed
entirely from retrained attackers. The ratio $\sum_i L/\sum_i M$ measures how much of the true
risk the scan recovers at a certification cost of $k$ retrained backgrounds per field instead
of all $\binom{|H|-1}{\le K}$.
